\documentclass[11pt]{article}
\usepackage[margin=1in]{geometry}
\usepackage{amsmath,amssymb,amsthm}
\usepackage{xurl}
\usepackage{hyperref}
\sloppy
\let\oldus\_ \renewcommand{\_}{\oldus\allowbreak}
\newtheorem{theorem}{Theorem}
\newtheorem{lemma}[theorem]{Lemma}
\theoremstyle{definition}
\newtheorem{definition}[theorem]{Definition}
\title{Conjecture 00000001543 is false:\\ $n$ collinear points on the irreducible curve $y=0$ determine only $n-1$ distances}
\author{Jackmeson1}
\date{2026-10-05}
\begin{document}
\maketitle

\section{The conjecture}
Conjecture 00000001543 (English text, verbatim): ``Definition: Distinct distances on algebraic
curves. Conjecture: Any n points on an irreducible real algebraic curve determine at least
$c\cdot n^{4/3}$ distinct distances (the curve distance exponent).'' The Chinese text says the
same: any $n$ points on any irreducible real algebraic curve determine at least $c\cdot n^{4/3}$
distinct distances.

\paragraph{Reading.} Points lie in the Euclidean plane $\mathbb R^2$ with the Euclidean distance.
An irreducible real algebraic curve is the zero set $Z(f)=\{p\in\mathbb R^2: f(p)=0\}$ of an
irreducible polynomial $f\in\mathbb R[x,y]$; we additionally require $Z(f)$ to be infinite, which
only shrinks the class (it excludes degenerate zero sets such as $Z(x^2+y^2)=\{0\}$). The number
of distinct distances of a finite set $P$ is $\#\{|p-q| : p,q\in P,\ p\neq q\}$. The constant $c$
is positive and independent of $n$. We refute the following readings, and claim no more:
\begin{itemize}
\item[(U)] there is one $c>0$ that works for every curve and every finite point set on it;
\item[(C)] every curve $C$ has its own $c_C>0$ and $N_C$ such that every set of $n\ge N_C$
points on $C$ determines at least $c_C n^{4/3}$ distances.
\end{itemize}
(C) is the weakest of the natural quantifier orders (constant depending on the curve, bound only
for large $n$), so (U) and every intermediate order (constant depending on the degree, bound for
all $n$) are refuted too.

\paragraph{Not refuted.} The statement as written does not exclude lines or circles. The known
theorem of Pach and de Zeeuw (Section~\ref{sec:rel}) does exclude them, and the reading
``irreducible curves other than lines and circles'' is a true theorem; we do not refute it.

\section{The counterexample}
\begin{lemma}\label{lem:curve}
The $x$-axis $L=\{(x,y): y=0\}=Z(y)$ is an irreducible real algebraic curve.
\end{lemma}
\begin{proof}
The variable $y$ is a prime, hence irreducible, element of $\mathbb R[x,y]$, and $Z(y)$ contains
the infinitely many points $(t,0)$, $t\in\mathbb R$.
\end{proof}

\begin{lemma}\label{lem:count}
Let $P_n=\{(k,0): k=0,1,\dots,n-1\}\subset L$. Then $|P_n|=n$ and $P_n$ determines at most $n-1$
distinct distances.
\end{lemma}
\begin{proof}
$|(a,0)-(b,0)|=|a-b|$, and for $a\neq b$ in $\{0,\dots,n-1\}$ this is an integer in
$\{1,\dots,n-1\}$.
\end{proof}

\begin{lemma}\label{lem:growth}
For every $c>0$ and every integer $n\ge \lceil c^{-3}\rceil+1$ we have $n-1<c\,n^{4/3}$.
\end{lemma}
\begin{proof}
$n>c^{-3}$ gives $(c\,n^{1/3})^3=c^3n>1$, so $c\,n^{1/3}>1$ and
$n-1<n<n\cdot c\,n^{1/3}=c\,n^{4/3}$.
\end{proof}

\begin{theorem}\label{thm:main}
For every $c>0$ there is $N$ such that for every $n\ge N$ some $n$ points of the irreducible curve
$L$ determine fewer than $c\,n^{4/3}$ distinct distances. Consequently readings (U) and (C) are
false.
\end{theorem}
\begin{proof}
Take $N=\lceil c^{-3}\rceil+1$ and $P_n$; combine Lemmas~\ref{lem:count} and~\ref{lem:growth}.
For (U): apply this to the given $c$ and $L$. For (C): given $c_L,N_L$, apply it with $c=c_L$
and $n=\max(N_L,N)$.
\end{proof}
The same failure occurs on circles (a regular $n$-gon determines $\lfloor n/2\rfloor$ distances);
we only use the line, and the circle remark is not formalized.

\section{Lean formalization}
File \texttt{lean/Conjecture1543/Basic.lean} (173 lines, only \texttt{import Mathlib}).
\begin{itemize}
\item \texttt{Plane := EuclideanSpace R (Fin 2)} (Euclidean metric; coordinates \texttt{p 0}, \texttt{p 1}).
\item \texttt{zeroSet f}: $\{p : \texttt{MvPolynomial.eval}\,(i\mapsto p\,i)\,f=0\}$ for
\texttt{f : MvPolynomial (Fin 2) R}.
\item \texttt{IsIrreducibleRealCurve C}: $\exists f$, \texttt{Irreducible f} $\wedge$
\texttt{C = zeroSet f} $\wedge$ \texttt{C.Infinite}.
\item \texttt{distances P}: the image of $(p,q)\mapsto \mathrm{dist}\,p\,q$ over
\texttt{P.offDiag} (ordered pairs of distinct points); its \texttt{card} is the number of
distinct distances.
\item \texttt{xAxis := zeroSet (X 1)}; \texttt{xAxis\_isIrreducibleRealCurve} (Lemma~\ref{lem:curve},
via Mathlib's \texttt{MvPolynomial.X\_prime}).
\item \texttt{collinear n}, \texttt{collinear\_card}, \texttt{collinear\_subset},
\texttt{distances\_collinear\_card\_le} (card $\le n-1$; Lemma~\ref{lem:count}).
\item \texttt{eventually\_lt} (Lemma~\ref{lem:growth}; real power \texttt{Real.rpow}).
\item \texttt{xAxis\_violates} (Theorem~\ref{thm:main}),
\texttt{not\_uniform} (reading U), \texttt{not\_per\_curve\_eventually} (reading C).
\end{itemize}
Axiom audit: \texttt{\#print axioms} for \texttt{xAxis\_violates}, \texttt{not\_uniform} and
\texttt{not\_per\_curve\_eventually} each reports \texttt{[propext, Classical.choice, Quot.sound]}.
No \texttt{sorry}, no \texttt{native\_decide}, no added axioms.

\section{Relation to the known theorem and to the earlier submission}\label{sec:rel}
Pach and de Zeeuw, \emph{Distinct distances on algebraic curves in the plane},
\url{https://arxiv.org/abs/1308.0177}, abstract: ``We prove that the number of distinct distances
determined by $P$ is at least $c_d n^{4/3}$, unless $C$ contains a line or a circle.'' The
conjecture's text drops the exception ``unless $C$ contains a line or a circle'', and the
exception is necessary: Theorem~\ref{thm:main} is exactly the line case.

\paragraph{Relation to the earlier submission.} PR \#221 (by orionsheep) used the same
counterexample, but its Lean theorem \texttt{no\_constant} was pure natural-number arithmetic. The
review rejected it because ``the actual counterexample (collinear points determine $n-1$
distances) is never formalized --- curves, points, and distances appear only in prose/Python.''
This submission formalizes the objects themselves: the Euclidean plane, irreducible polynomials in
\texttt{MvPolynomial (Fin 2) R} and their zero sets, finite point sets on the curve, and the
distinct-distance count as the cardinality of the image of the distance map. It proves the
violation for every $c>0$ and all large $n$, under both a uniform and a curve-dependent constant.
\end{document}
